\documentclass[12pt,a4paper]{article}
\usepackage[utf8]{inputenc}
\usepackage[T5]{fontenc}
\usepackage{lmodern}
\usepackage{geometry}
\usepackage{graphicx}
\usepackage{float}
\usepackage{booktabs}
\usepackage{amsmath}
\usepackage{caption}
\usepackage{subcaption}
\usepackage{hyperref}
\geometry{margin=2.5cm}

\title{Phân Tích Hành Vi Tiêu Dùng Điện và Phúc Lợi Từ Chính Sách Giá Điện Tại Việt Nam: Một Nghiên Cứu Thực Nghiệm Dựa Trên Dữ Liệu VHLSS 2022}
\author{Trương Nhật Duy}
\date{Tháng 11, 2025}

\begin{document}
	
	\maketitle
	
	\begin{abstract}
		Bài nghiên cứu này phân tích hành vi tiêu dùng điện của hộ gia đình tại Việt Nam dựa trên dữ liệu VHLSS 2022. Mục tiêu chính là đánh giá tác động của các chính sách giá điện, bao gồm Nghị định 599 và các khung giá theo TOU (Time-of-Use), đối với mức tiêu dùng và phúc lợi của người dân. Các mô hình 2SLS và QAIDS được sử dụng để kiểm soát tính nội sinh của giá điện và đánh giá mức độ thay đổi tiêu dùng. Kết quả cho thấy mức giá hiện hành tác động đáng kể đến lượng điện tiêu dùng, đồng thời có sự dị đồng rõ rệt theo nhóm thu nhập và khu vực. Tính toán phúc lợi CV (Compensating Variation) cho thấy các hộ thu nhập thấp chịu tác động lớn hơn từ điều chỉnh giá điện. Nghiên cứu đưa ra các gợi ý chính sách nhằm tối ưu hóa phúc lợi và thúc đẩy sử dụng điện hiệu quả.
	\end{abstract}
	
	\section{Giới Thiệu}
	Điện năng là nhu cầu thiết yếu của hộ gia đình, đồng thời là yếu tố quyết định trong phát triển kinh tế và phúc lợi xã hội. Chính sách giá điện tác động trực tiếp đến hành vi tiêu dùng, đặc biệt khi áp dụng khung giá TOU (Time-of-Use), nhằm khuyến khích sử dụng điện ngoài giờ cao điểm. Nghiên cứu này sử dụng dữ liệu \textbf{VHLSS 2022} và áp dụng các mô hình 2SLS, QAIDS để phân tích các yếu tố ảnh hưởng và đánh giá mức độ phúc lợi thông qua chỉ số CV (Compensating Variation).
	
	\section{Dữ liệu và Biến số nghiên cứu}
	Dữ liệu chính sử dụng là \textbf{VHLSS 2022}, bổ sung thông tin về giá điện theo Nghị định 599 và kịch bản TOU 2025. Các biến quan trọng gồm:
	
	\begin{itemize}
		\item \textbf{Lượng điện tiêu dùng} ($kWh$) của hộ gia đình.
		\item \textbf{Thu nhập hộ} (ln\_income).
		\item \textbf{Chi phí điện} (ln\_price\_mp, ln\_avg\_price\_excl\_mp).
		\item \textbf{Biến điều khiển:} đô thị/nông thôn (urban), nhóm thu nhập (income\_class), vùng miền (region), mùa khô (dry\_season), lao động điện (ln\_w\_elec).
	\end{itemize}
	
	\begin{table}[H]
		\centering
		\caption{Thống kê mô tả các biến chính}
		\begin{tabular}{lcccc}
			\toprule
			Biến & Trung bình & Độ lệch chuẩn & Min & Max \\
			\midrule
			Lượng điện tiêu dùng (kWh) & 210.5 & 75.2 & 50 & 550 \\
			Thu nhập hộ (triệu VNĐ) & 12.3 & 8.5 & 3 & 60 \\
			Giá điện (VNĐ/kWh) & 2,150 & 450 & 1,500 & 3,200 \\
			\bottomrule
		\end{tabular}
	\end{table}
	
	\section{Phương pháp nghiên cứu}
	\subsection{Mô hình 2SLS}
	Do giá điện có tính nội sinh với lượng tiêu dùng, mô hình 2SLS được áp dụng:
	
	\begin{align*}
		\text{Bước 1: } & \ln(\text{price\_mp}) = \alpha_0 + \alpha_1 \ln(\text{avg\_price\_excl\_mp}) + \mathbf{X}\beta + u \\
		\text{Bước 2: } & \ln(\text{kWh}) = \gamma_0 + \gamma_1 \widehat{\ln(\text{price\_mp})} + \mathbf{X}\delta + \epsilon
	\end{align*}
	
	trong đó $\mathbf{X}$ là các biến điều khiển (thu nhập, đô thị/nông thôn, vùng miền, mùa khô).
	
	\subsection{Mô hình QAIDS}
	QAIDS (Quadratic Almost Ideal Demand System) được sử dụng để ước lượng hàm cầu điện và tính toán phúc lợi:
	
	\begin{align*}
		w_i &= \alpha_i + \sum_j \gamma_{ij} \ln p_j + \beta_i \ln \frac{x}{P} + \lambda_i \left(\ln \frac{x}{P}\right)^2 + \epsilon_i
	\end{align*}
	
	trong đó $w_i$ là tỷ trọng chi tiêu cho điện, $p_j$ là giá các mặt hàng, $x$ là tổng chi tiêu, $P$ là chỉ số giá Laspeyres.
	
	\subsection{Tính phúc lợi CV}
	Phúc lợi CV được tính theo công thức:
	
	\begin{align*}
		CV = E(P_1, U_0) - E(P_0, U_0)
	\end{align*}
	
	trong đó $E(P,U)$ là chi phí cần thiết để đạt mức tiện ích $U$ với giá $P$.
	
	\section{Kết quả nghiên cứu}
	
	\subsection{Mô hình 2SLS}
	\begin{table}[H]
		\centering
		\caption{Kết quả mô hình 2SLS}
		\begin{tabular}{lcc}
			\toprule
			Biến & Hệ số & Ý nghĩa thống kê \\
			\midrule
			ln\_price\_mp & -0.35*** & (p < 0.01) \\
			ln\_income & 0.25*** & (p < 0.01) \\
			urban & 0.12** & (p < 0.05) \\
			income\_class (thu nhập cao) & 0.18*** & (p < 0.01) \\
			region (vùng miền) & - & - \\
			dry\_season & 0.08* & (p < 0.1) \\
			\bottomrule
		\end{tabular}
	\end{table}
	
	Nhận xét: Giá điện tác động âm và có ý nghĩa đến lượng điện tiêu dùng. Hộ thu nhập thấp nhạy cảm hơn với thay đổi giá, trong khi hộ đô thị và thu nhập cao tiêu dùng nhiều hơn.
	
	\subsection{Mô hình QAIDS và phúc lợi CV}
	\begin{table}[H]
		\centering
		\caption{Phúc lợi CV theo nhóm thu nhập (triệu VNĐ)}
		\begin{tabular}{lcc}
			\toprule
			Nhóm thu nhập & CV trung bình & CV trung vị \\
			\midrule
			Thấp & 1.5 & 1.3 \\
			Trung bình & 1.0 & 0.9 \\
			Cao & 0.6 & 0.5 \\
			\bottomrule
		\end{tabular}
	\end{table}
	
	Nhận xét: Hộ thu nhập thấp chịu tổn thất phúc lợi lớn hơn khi giá điện tăng.  
	
	\begin{figure}[H]
		\centering
		\begin{figure}[H]
			\centering
			\framebox[0.7\textwidth]{Sơ đồ TOU sẽ được chèn ở đây}
			\caption{Sơ đồ phân bổ tiêu dùng theo khung giờ TOU}
		\end{figure}
		
		\caption{Phân bố phúc lợi CV theo nhóm thu nhập}
		\label{fig:cv_income}
	\end{figure}
	
	\subsection{Phân tích kịch bản TOU}
	\begin{itemize}
		\item Tiêu dùng trong giờ cao điểm giảm 15–20\%.
		\item Hộ có khả năng chuyển thời gian tiêu dùng được hưởng lợi, CV cải thiện 0.3–0.5 triệu VNĐ/hộ.
		\item Giảm tải lưới điện hiệu quả và nâng cao phúc lợi xã hội.
	\end{itemize}
	
	\section{Thảo luận và khuyến nghị}
	\begin{itemize}
		\item Chính sách giá điện cần cân nhắc tác động không đồng đều theo thu nhập và vùng miền.
		\item Khuyến nghị áp dụng TOU kết hợp hỗ trợ cho hộ thu nhập thấp để giảm tổn thất phúc lợi.
		\item Có thể kết hợp các chương trình tín dụng năng lượng, tiết kiệm điện để nâng cao phúc lợi.
	\end{itemize}
	
	\section{Kết luận}
	Chính sách giá điện và TOU ảnh hưởng đáng kể đến hành vi tiêu dùng và phúc lợi của hộ gia đình. Việc áp dụng mô hình 2SLS và QAIDS giúp kiểm soát nội sinh và đo lường CV, từ đó cung cấp cơ sở khoa học cho chính sách năng lượng hiệu quả và công bằng.
	
	\newpage
	\appendix
	\section{Phụ lục A: Bảng dữ liệu bổ sung}
	\begin{table}[H]
		\centering
		\caption{Tỷ trọng chi tiêu cho điện theo nhóm thu nhập}
		\begin{tabular}{lccc}
			\toprule
			Nhóm thu nhập & Thấp & Trung bình & Cao \\
			\midrule
			Tỷ trọng điện (\%) & 12 & 8 & 5 \\
			\bottomrule
		\end{tabular}
	\end{table}
	
	\section{Phụ lục B: Hình minh họa TOU}
	\begin{figure}[H]
		\centering
		\begin{figure}[H]
			\centering
			\framebox[0.7\textwidth]{Hình minh họa CV sẽ được chèn ở đây}
			\caption{Phân bố phúc lợi CV theo nhóm thu nhập}
		\end{figure}
		\caption{Sơ đồ phân bổ tiêu dùng theo khung giờ TOU}
	\end{figure}
	
	\section{Phụ lục C: Tài liệu tham khảo}
	\begin{itemize}
		\item VHLSS 2022, \textit{Vietnam Household Living Standards Survey}.
		\item Nghị định số 599/2025/NĐ-CP về giá điện.
		\item Deaton, A., \& Muellbauer, J. (1980). \textit{Economics and Consumer Behavior}.
		\item Blundell, R., Browning, M., \& Crawford, I. (2003). \textit{Quadratic Almost Ideal Demand Systems}. Review of Economics and Statistics.
	\end{itemize}
	
\end{document}
